\documentclass[aspectratio=169,10pt]{beamer}

% ============================================================================
%  FICTIONAL DEMO DECK — For Teaser / Promotional Use Only
%  All content is fabricated. No real people, papers, or institutions.
% ============================================================================

\usepackage[academic]{template-lib/template-lib}
\uselayout{text}
\uselayout{eq}
\uselayout{twocol}
\uselayout{1img}

\TLsetfoot{Demo Deck · Beamer Deck Auto}

\title{Sample Research Presentation}
\subtitle{A Fictional Example for Template Demonstration}
\author{A. N. Author}
\institute{Example University\\Department of Demonstrations}
\date{May 2026}

\begin{document}

% ── Title Slide ─────────────────────────────────────────────────────────────
\TLtitlecenter
  {Sample Research Presentation}
  {Template Demonstration}
  {A. N. Author}
  {Example University · May 2026}

% ── S1: Overview ────────────────────────────────────────────────────────────
\TLsection{Overview}

\begin{frame}{Project Goals}
  \begin{TLtext}
    \begin{itemize}
      \item \textbf{Goal 1:} Demonstrate the three-tier template system
      \item \textbf{Goal 2:} Show layout flexibility across content types
      \item \textbf{Goal 3:} Illustrate component reuse with consistent styling
    \end{itemize}
  \end{TLtext}

  \vspace{0.5em}
  \TLtakeaway{%
    This deck contains no real data, no real people, and no real institutions.%
  }
\end{frame}

% ── S2: Method ──────────────────────────────────────────────────────────────
\TLsection{Method}

\begin{frame}{Mathematical Framework}
  \TLeqsingle{%
    \mathcal{L}(\theta) = \sum_{i=1}^{n} \ell\bigl(f_\theta(x_i), y_i\bigr) + \lambda R(\theta)
  }{%
    Standard regularized empirical risk minimization.%
  }

  \vspace{0.5em}
  \begin{TLtext}
    \begin{itemize}
      \item $f_\theta$: parametric model family
      \item $\ell$: per-sample loss function
      \item $R(\theta)$: regularizer (e.g., $\|\theta\|_2^2$)
    \end{itemize}
  \end{TLtext}
\end{frame}

\begin{frame}{Two-Column Comparison}
  \begin{columns}[T]
    \begin{column}{0.48\textwidth}
      \TLinfoblock{Approach A}{%
        \begin{itemize}
          \item Simple implementation
          \item Fast convergence
          \item Lower memory
        \end{itemize}%
      }
    \end{column}
    \begin{column}{0.48\textwidth}
      \TLinfoblock{Approach B}{%
        \begin{itemize}
          \item Higher accuracy
          \item Better generalization
          \item Scales well
        \end{itemize}%
      }
    \end{column}
  \end{columns}

  \vspace{0.5em}
  \TLtakeaway{%
    Trade-offs depend on dataset size and latency constraints.%
  }
\end{frame}

% ── S3: Results ─────────────────────────────────────────────────────────────
\TLsection{Results}

\begin{frame}{Key Findings}
  \begin{TLtext}
    \begin{itemize}
      \item Method achieves \textbf{94.2\%} accuracy on benchmark
      \item Training time reduced by \textbf{35\%} vs. baseline
      \item Model size: \textbf{12M} parameters
    \end{itemize}
  \end{TLtext}

  \vspace{0.3em}
  \TLresultblock{Result}{%
    The proposed approach outperforms all baselines while maintaining
    computational efficiency suitable for edge deployment.%
  }
\end{frame}

% ── S4: Conclusion ──────────────────────────────────────────────────────────
\TLsection{Conclusion}

\begin{frame}{Summary}
  \begin{TLtext}
    \begin{enumerate}
      \item Three-tier template system enables rapid deck construction
      \item Layout optimizer selects appropriate page structure automatically
      \item Design grammar checker ensures visual quality before presentation
    \end{enumerate}
  \end{TLtext}

  \vspace{0.5em}
  \TLtakeaway{%
    Beamer Deck Auto: systematic slide design for researchers.%
  }
\end{frame}

\end{document}
